% !PRINTED 409
\weberpage{409}
Wenn die $n^{\mathrm{te}}$ Potenz einer Zahl $r$ gleich 1 ist, wenn also die Gleichung

\weq{1}{r^n = 1}
für ein ganzes positives $n$ befriedigt ist, so heisst $r$ eine $n^{\mathrm{te}}$
\gesperrt{Einheitswurzel} oder eine \gesperrt{Einheitswurzel vom Grade $n$}. Es sind also alle $n^{\mathrm{ten}}$
Einheitswurzeln, und nur diese, Wurzeln der Gleichung $n^{\mathrm{ten}}$ Grades

\weq{2}{f(x) = x^n - 1 = 0.}
Es ist

\weq{3}{f'(x) = n\,x^{n-1},}
und folglich hat $f(x)$ mit $f'(x)$ keinen Theiler gemein; also hat $f(x)$ keine
mehrfachen Wurzeln, und es giebt $n$ und nicht mehr von einander verschiedene
$n^{\mathrm{te}}$ Einheitswurzeln.

Ist $r$ eine $n^{\mathrm{te}}$ Einheitswurzel, so ist es auch jede ganze Potenz von $r$,
denn aus $r^n = 1$ folgt $r^{kn} = 1$, wenn $k$ eine beliebige positive oder negative
ganze Zahl ist (auch $k = 0$ nicht ausgeschlossen); also ist auch $r^k$ eine $n^{\mathrm{te}}$
Einheitswurzel. Da es aber nur $n$ Einheitswurzeln vom Grade $n$ giebt, so sind die
Potenzen $r^k$ nicht alle von einander verschieden. Hierüber gilt nun Folgendes:

Wenn zwei Zahlen $k, k'$ sich um ein Vielfaches von $n$ unterscheiden, wenn also

\weq{4}{k' \equiv k \pmod{n}}
ist, so ist auch

\weq{5}{r^k = r^{k'}.}
Denn ist $k' = k + h\,n$, so ist

\[
r^{k'} = r^k r^{h n},
\]
woraus, da $r^n = 1$ ist, die Gleichung (5) folgt.

Es sind also in der Reihe der Zahlen

\weq{6}{1, r, r^2, r^3 \ldots r^{n-1}}
gewiss alle von einander verschiedenen $r^k$ enthalten; aber es müssen nicht umgekehrt
die Gr\"ossen (6) alle von einander verschieden sein.

Nehmen wir an, es seien $k$ und $k' = k + \mu$ zwei Zahlen der Reihe

\[
0, 1, 2 \ldots n-1
\]
und

\[
r^k = r^{k+\mu},
\]
%CONT>>
